% ============================================================================
% Appendix: Poisson Ratio Determination (ASTM D638 Type I, Tough1500)
%
% Required packages (add to preamble if not already present):
%   \usepackage{amsmath}
%   \usepackage{amssymb}
%   \usepackage{booktabs}     % for the comparison tables
%   \usepackage{siunitx}      % optional, used for unit formatting
% ============================================================================

\chapter{Poisson Ratio Determination}
\label{app:poisson-ratio}

This appendix documents the measurement of Poisson's ratio for the Tough1500
photopolymer used throughout the AM\_Creep\_Analysis pipeline (TSSP master curve
construction, edge-tracking analysis, and FEM input). The result replaces the
nominal placeholder $\nu = 0.43$ that had been used in earlier processing.
The analysis script is \texttt{ASRMD638TypeI/analyze\_dogbone.py}, operating on
ASTM D638 Type~I tensile data in \texttt{ASRMD638TypeI/\{0,45\}/*.csv} with
specimen dimensions tabulated in \texttt{ASTM~Type~I~-~Sheet1.csv}; outputs are
written to \texttt{ASRMD638TypeI/out/dogbone\_results.csv} and accompanying
PNGs.

\section{Result}
\label{sec:poisson-result}

The measured Poisson ratio for the $0^{\circ}$ print orientation, averaged over
five specimens (\textit{p-5} through \textit{p-9}), is
%
\begin{equation}
  \nu \;=\; 0.421 \pm 0.025 \quad (n = 5),
  \label{eq:poisson-result}
\end{equation}
%
with per-sample values $\{0.456,\,0.418,\,0.411,\,0.433,\,0.390\}$ and
$r^{2}_{\nu} \geq 0.9994$ on every fit. The $45^{\circ}$ specimens were
processed but not used: subsequent tests had moved to a video extensometer with
a cam gauge length of approximately $59$--$68\,\mathrm{mm}$, which exceeds the
Type~I parallel section length ($\sim 50\,\mathrm{mm}$) and therefore bridges
into the fillet region.

\section{Correction to the Original Formula}
\label{sec:poisson-formula}

The original \texttt{fit\_poisson()} routine regressed the raw transverse
displacement against the raw axial displacement and reported the slope as
$\nu$. This is incorrect when the two extensometers reference different gauge
lengths, as is the case here.

The transverse clip-on extensometer has a $5\,\mathrm{mm}$ arm span, with its
blades resting on opposite specimen edges. The raw transverse channel is
therefore the true specimen width change $\delta w$, but the
\texttt{Ext.2(Strain)} channel reports $\delta w / 5\,\mathrm{mm}$, not
$\delta w / w_{\text{specimen}}$. The true transverse strain is obtained by
dividing $\delta w$ by the actual specimen width
($w_{\text{specimen}} \approx 13\,\mathrm{mm}$ for Type~I), yielding
%
\begin{equation}
  \nu \;=\; -\,\frac{\varepsilon_{\text{trans}}}{\varepsilon_{\text{axial}}}
        \;=\; \frac{L_{\text{axial gauge}}}{w_{\text{specimen}}}
        \,\cdot\, \mathrm{slope}\!\left(\Delta\,\text{trans}_{\text{disp}},\;
        \Delta\,\text{axial}_{\text{disp}}\right).
  \label{eq:poisson-correct}
\end{equation}
%
For the Type~I $0^{\circ}$ geometry with $L_{\text{axial gauge}} = 25\,\mathrm{mm}$
(Ext.1) and $w_{\text{specimen}} \approx 13\,\mathrm{mm}$, the correction factor
is approximately $1.92\times$. The legacy values of $0.21$--$0.26$ were biased
low by exactly that factor.

\section{Fit Window: Axial Strain, Not Stroke}
\label{sec:poisson-window}

Two stroke-based fit windows were trialled before the final strain-based window
was adopted.

\paragraph{Stroke $\leq 3\,\mathrm{mm}$.}
The specimen has already passed its proportional limit
(stress $\approx 30\,\mathrm{MPa}$ at axial strain $\approx 1.8\,\%$), so the
slope is biased low by plastic deformation.

\paragraph{Stroke $\leq 1\,\mathrm{mm}$.}
The regression has a non-zero intercept ($\approx +0.0016\,\mathrm{mm}$) from
grip take-up and clip-on seating drift. Over a small displacement window this
offset dominates and inflates the slope, producing unphysical $\nu > 0.5$ on
several specimens.

\paragraph{Final window: axial strain $0.1\,\%$--$0.5\,\%$ (via Ext.1).}
Using the axial extensometer's own strain channel skips grip take-up
(the extensometer zeroes itself once contact is established), and the upper
bound at $0.5\,\%$ strain ($\sim 10\,\mathrm{MPa}$ stress) stays well below
yield. With this window $r^{2}_{\nu} \geq 0.998$ on every retained sample.

\section{Sample Exclusion: Warm-up Drift}
\label{sec:poisson-exclusion}

Even with the cleaned strain window, the first three tested specimens returned
$\nu > 0.5$, as shown in Table~\ref{tab:poisson-perspecimen}.

\begin{table}[!ht]
  \centering
  \caption{Per-specimen Poisson ratio in the $0.1$--$0.5\,\%$ axial strain
           window. The break between \textit{p-4} and \textit{p-5} coincides
           exactly with the tested-order transition.}
  \label{tab:poisson-perspecimen}
  \begin{tabular}{lc}
    \toprule
    Sample & $\nu$ ($0.1$--$0.5\,\%$ window) \\
    \midrule
    p-1 & 0.578 \\
    p-3 & 0.588 \\
    p-4 & 0.539 \\
    \midrule
    p-5 & 0.456 \\
    p-6 & 0.418 \\
    p-7 & 0.411 \\
    p-8 & 0.433 \\
    p-9 & 0.390 \\
    \bottomrule
  \end{tabular}
\end{table}

The discontinuity falls exactly at the boundary between the early and later
specimens of the testing session and is attributed to warm-up plus transverse
clip-on seating drift on the first few tests. These specimens are excluded
via the constant \verb|EXCLUDE_SAMPLES = {0: {1, 3, 4}}| in
\texttt{analyze\_dogbone.py}.

\section{Type~I vs.\ Type~V Comparison}
\label{sec:poisson-typecompare}

\begin{table}[!ht]
  \centering
  \caption{Tensile metrics for Type~I ($0^{\circ}$, $n = 5$) and Type~V
           ($n = 21$) specimens. Modulus values are not directly comparable;
           see text.}
  \label{tab:poisson-typeIvsV}
  \begin{tabular}{lcc}
    \toprule
    Metric & Type~I ($0^{\circ}$, $n=5$) & Type~V ($n=21$) \\
    \midrule
    $E$    & $1902 \pm 62\,\mathrm{MPa}$ (Ext.1, $25\,\mathrm{mm}$ GL)
           & $347 \pm 21\,\mathrm{MPa}$ (apparent, from stroke) \\
    UTS    & $46.8 \pm 0.4\,\mathrm{MPa}$ & $56.9 \pm 1.8\,\mathrm{MPa}$ \\
    Area   & $42\,\mathrm{mm}^{2}$ & $12.8\,\mathrm{mm}^{2}$ \\
    \bottomrule
  \end{tabular}
\end{table}

The two modulus columns in Table~\ref{tab:poisson-typeIvsV} are not
comparable. Type~V strains were computed from crosshead stroke, which contains
apparatus compliance dominating the true gauge deformation by roughly
$6\times$. Type~I data with Ext.1 returns the true material modulus, in good
agreement with the Tough1500 datasheet value of $\sim 2.0\,\mathrm{GPa}$.
Type~V tensile data are excluded from all subsequent analysis because no
extensometer record is available for those tests.

The UTS gap is largely physical: a Weibull size effect (the Type~V
cross-section is approximately $3\times$ smaller) combined with an
effectively higher strain rate on the smaller gauge accounts for the
difference.

\section{Apparatus Compliance (Side Finding)}
\label{sec:poisson-apparatus}

A by-product of the simultaneous stroke and Ext.1 measurement is a direct
estimate of the apparatus compliance. Fitting $(\text{stroke} - \text{axial})$
versus force in the $0.05$--$0.25\,\%$ elastic window over the five Type~I
$0^{\circ}$ specimens gives
%
\begin{equation}
  C_{\text{apparatus}} \;=\; 1.84 \pm 0.04\;\mathrm{\mu m/N}
  \qquad (n = 5,\; \mathrm{CV} = 2\,\%),
  \label{eq:apparatus-compliance}
\end{equation}
%
which decomposes (using the gauge stiffness inferred from Ext.1) as
%
\begin{align}
  C_{\text{gauge}}     &\approx 0.31\;\mathrm{\mu m/N}, \\
  C_{\text{apparatus}} &\approx 1.84\;\mathrm{\mu m/N}, \\
  C_{\text{total}}     &\approx 2.15\;\mathrm{\mu m/N}.
\end{align}
%
The apparatus compliance lumps together machine compliance
(frame, load cell, grips) and the non-gauge portion of the specimen (the
parallel section outside Ext.1, plus fillets and tabs). An analytical
breakdown suggests roughly $1.0\,\mathrm{\mu m/N}$ of machine compliance plus
$0.8\,\mathrm{\mu m/N}$ of non-gauge specimen compliance, but only the lumped
quantity is what would be subtracted from Type~I stroke-only data as a
correction. A clean separation would require either a stiff reference bar or
a multi-gauge-length test campaign and is deferred; a follow-on script
\texttt{scripts/tensile/calculate\_apparatus\_compliance.py} is planned to
verify load-independence of $C_{\text{apparatus}}$ across varied peak loads.

\section{Summary}
\label{sec:poisson-summary}

The Tough1500 Poisson ratio used as input to the TSSP, edge-tracking, and FEM
pipelines is
%
\[
  \boxed{\,\nu = 0.421 \pm 0.025\,}
\]
%
measured on $n = 5$ ASTM D638 Type~I specimens at $0^{\circ}$ print
orientation, fit over a $0.1$--$0.5\,\%$ axial strain window with the
gauge-length correction of equation~\eqref{eq:poisson-correct} applied. This
supersedes the previously used placeholder $\nu = 0.43$.
